% --- les trois étiquettes de provenance -------------------------------------------------------
% La prose porte les mêmes étiquettes que les colonnes de l'entrepôt, et pour la même raison :
% un lecteur doit pouvoir dire, sur n'importe quelle affirmation, d'où elle vient.
%
%   DOC  une entrée du registre des sources l'étaye. Elle se cite par \cite{...}, et la clé
%        renvoie au fichier bibliographique produit à partir du registre.
%   OBS  un relevé d'observation l'étaye. Aucune source publiée ne la porte : elle a été vue.
%   HYP  aucune des deux. C'est une convention posée, et le rapport la déclare comme telle.
%
% Chaque fichier de chapitre déclare en tête l'ensemble de ce sur quoi il repose, sous ces trois
% étiquettes, et un contrôle vérifie que la déclaration coïncide exactement avec ce que le fichier
% porte — dans les deux sens.

% --- ces deux étiquettes NE COMPOSENT PLUS RIEN ------------------------------------------------
% Elles portaient un exposant `[obs.]` et `[conv.]` au fil de la prose. Elles ne l'impriment plus,
% et le motif est le DESTINATAIRE du document : ce rapport s'adresse à un encadrant de stage, non
% à une relecture adverse. Un exposant tous les deux paragraphes est un appareil de défense, et il
% se lit comme tel.
%
% L'APPAREIL RESTE ENTIÈREMENT CÂBLÉ, IL CESSE SEULEMENT D'IMPRIMER. L'argument reste l'identifiant
% du relevé ou de la convention ; chaque fichier de chapitre continue de déclarer en tête ce sur
% quoi il repose ; et `tests/test_provenance_des_chapitres.py` continue de confronter la
% déclaration à ce que le fichier porte, dans les deux sens. CE CONTRÔLE LIT LES SOURCES DE
% COMPOSITION, JAMAIS LE PDF : rien de ce qu'il vérifie ne dépend de ce que ces deux commandes
% composent, et les rendre muettes ne l'affaiblit d'aucune façon. Retirer un `\releve{...}` d'une
% phrase reste rouge, exactement comme avant.
%
% Les citations `\cite{...}` de l'étiquette DOC, elles, restent visibles : renvoyer à une source
% publiée est l'usage académique normal, et ce n'est pas un appareil de défense.

% Un relevé d'observation. L'argument est son identifiant, court et stable ; il n'est pas composé.
\newcommand{\releve}[1]{}

% Une convention posée. Même forme, même règle.
\newcommand{\convention}[1]{}

% --- ce qui reste à rédiger --------------------------------------------------------------------
% Un paragraphe qui relève de la rédaction personnelle et que le squelette ne peut pas écrire.
% Le premier argument est son identifiant d'observation, le second ce qu'il portera.
%
% La marque est VISIBLE à la compilation, toujours : un rapport remis en la portant se voit à
% l'œil. Le contrôle, lui, s'appuie sur l'état déclaré du document — il rougit en état de remise
% et s'abstient en état de brouillon, où un paragraphe non encore écrit est normal.
\newcommand{\aRediger}[2]{%
  \par\medskip\noindent\fbox{%
    \begin{minipage}{0.95\linewidth}%
      \textbf{À rédiger} — #2%
    \end{minipage}%
  }\par\medskip%
}
